\begin{center}
\hypertarget{native-043-article027}{}\pdfbookmark[2]{27.}{article-027}27.
\end{center}
Si aequatio IV ita differentiatur, ut $u$, $v$, $\psi$ simul ut variabiles tractentur, prodit
\[
\begin{aligned}
\frac{\mathrm{d}u}{u}&=\frac{\sin\psi\,\mathrm{d}v+\sin v\,\mathrm{d}\psi}{2\cos\frac{1}{2}(v-\psi)\cos\frac{1}{2}(v+\psi)}\\
&=\frac{r\operatorname{tang}\psi}{p}\,\mathrm{d}v+\frac{r\sin v}{p\cos\psi}\,\mathrm{d}\psi.
\end{aligned}
\]

Differentiando perinde aequationem XI, inter variationes differentiales quantitatum $u$, $\psi$, $N$ emergit relatio
\SourceLeaf{044}
\[
\frac{\mathrm{d}N}{\lambda}=\left\{\frac{1}{2}e\left(1+\frac{1}{uu}\right)-\frac{1}{u}\right\}\mathrm{d}u+\frac{(uu-1)\sin\psi}{2u\cos\psi^2}\,\mathrm{d}\psi
\]
sive
\[
\frac{\mathrm{d}N}{\lambda}=\frac{r}{bu}\,\mathrm{d}u+\frac{r\sin v}{b\cos\psi}\,\mathrm{d}\psi.
\]
Hinc eliminando $\mathrm{d}u$ adiumento aequationis praecedentis obtinemus
\[
\frac{\mathrm{d}N}{\lambda}=\frac{rr}{bb\operatorname{tang}\psi}\,\mathrm{d}v+\left(1+\frac{r}{p}\right)\frac{r\sin v}{b\cos\psi}\,\mathrm{d}\psi
\]
sive
\[
\begin{aligned}
\mathrm{d}v&=\frac{bb\operatorname{tang}\psi}{\lambda rr}\,\mathrm{d}N-\left(\frac{b}{r}+\frac{b}{p}\right)\frac{\sin v\operatorname{tang}\psi}{\cos\psi}\raisebox{-.6ex}{.}\,\mathrm{d}\psi\\
&=\frac{bb\operatorname{tang}\psi}{\lambda rr}\,\mathrm{d}N-\left(1+\frac{p}{r}\right)\frac{\sin v}{\sin\psi}\,\mathrm{d}\psi.
\end{aligned}
\]

\begin{center}
\hypertarget{native-044-article028}{}\pdfbookmark[2]{28.}{article-028}28.
\end{center}
Differentiando aequationem X, omnibus $r$, $b$, $e$, $u$ pro variabilibus habitis, substituendo $\mathrm{d}e=\frac{\sin\psi}{\cos\psi^2}\,\mathrm{d}\psi$, eliminandoque $\mathrm{d}u$ adiumento aequationis inter $\mathrm{d}N$, $\mathrm{d}u$, $\mathrm{d}\psi$ in art. praec. traditae, prodit
\[
\begin{aligned}
\mathrm{d}r&=\frac{r}{b}\,\mathrm{d}b+\frac{bbe(uu-1)}{2\lambda ur}\,\mathrm{d}N\\
&\quad+\frac{b}{2\cos\psi^2}\left\{\left(u+\frac{1}{u}\right)\sin\psi-\left(u-\frac{1}{u}\right)\sin v\right\}\mathrm{d}\psi.
\end{aligned}
\]
Coëfficiens ipsius $\mathrm{d}N$ per aequ. VIII transit in $\frac{b\sin v}{\lambda\sin\psi}$; coëfficiens ipsius $\mathrm{d}\psi$ autem, substituendo, per aequ. IV, $u(\sin\psi-\sin v)=\sin(\psi-v)$, $\frac{1}{u}(\sin\psi+\sin v)=\sin(\psi+v)$, mutatur in $\frac{b\sin\psi\cos v}{\cos\psi^2}=\frac{p\cos v}{\sin\psi}$, ita ut habeatur
\[
\mathrm{d}r=\frac{r}{b}\,\mathrm{d}b+\frac{b\sin v}{\lambda\sin\psi}\,\mathrm{d}N+\frac{p\cos v}{\sin\psi}\,\mathrm{d}\psi.
\]
Quatenus porro $N$ ut functio ipsarum $b$ et $t$ spectatur, fit $\mathrm{d}N=\frac{N}{t}\,\mathrm{d}t-\frac{3}{2}\cdot\frac{N}{b}\,\mathrm{d}b$, quo valore substituto, $\mathrm{d}r$, ac perinde in art. praec. $\mathrm{d}v$, per $\mathrm{d}t$, $\mathrm{d}b$, $\mathrm{d}\psi$ expressae habebuntur. Ceterum quod supra monuimus etiam hic repetendum est, scilicet si angulorum $v$ et $\psi$ variationes non in partibus radii sed in secundis expressae concipiantur, vel omnes terminos qui $\mathrm{d}v$, $\mathrm{d}\psi$ continent per $206264{,}8$ dividi, vel omnes reliquos per hunc numerum multiplicari debere.
